% Lösungen Einheit 5 von 5 – Zusammengesetzte Körper und Anwendungen (zusammenbau v0.1)
\einheitenkopf{Einheit 5 von 5 · Zusammengesetzte Körper und Anwendungen}

%% TODO Lösung der Erklärzeile 25g) fehlt in der Bank
\erg{25}{a) Quader (Wohnteil) und Dreiecksprisma (Dach) \quad b) $8 \cdot 4 \cdot 2 + 8 \cdot 2 \cdot 2 = 64 + 32 = 96$ dm³ \quad c) $9 \cdot 3 \cdot 2 + 3 \cdot 3 \cdot 5 = 54 + 45 = 99$ cm³ \quad d) $7 \cdot 3 \cdot 3 + 4 \cdot 3 \cdot 2 = 63 + 24 = 87$ m³ \quad e) $5 \cdot 2 \cdot 1 + 5 \cdot 1 \cdot 2 = 10 + 10 = 20$ m³ \quad f) $9 \cdot 6 \cdot 1 + 3 \cdot 3 \cdot 7 = 54 + 63 = 117$ dm³ \quad g) \ldots \quad h) Quader $11 \cdot 7 \cdot 3 = 231$ m³; Dach $7 \cdot 3 : 2 \cdot 11 = 115{,}5$ m³; $231 + 115{,}5 = 346{,}5$ m³ \quad i) $8 \cdot 4 \cdot 1{,}4 = 44{,}8$ m³; $\tfrac{1}{2} \cdot \pi \cdot 2^2 \cdot 1{,}4 \approx 8{,}8$ m³; $44{,}8 + 8{,}8 = 53{,}6$ m³ \quad j) $(7 \cdot 4 + \tfrac{1}{2} \cdot \pi \cdot 2^2) \cdot 1{,}6 \approx 54{,}9$ m³ – die Tiefe gilt für Rechteck und Halbkreis, darum die Klammer \quad k) Karton $24 \cdot 24 \cdot 9 = 5\,184$ cm³; Torte $\pi \cdot 11^2 \cdot 8 \approx 3\,041{,}1$ cm³; $5\,184 - 3\,041{,}1 = 2\,142{,}9$ cm³ \quad l) $37$ cm × $37$ cm × $56$ cm, weil der Durchmesser $= 2 \cdot 18 = 36$ cm in Länge und Breite passen muss und die Höhe $55$ cm \quad m) $4 \cdot 4 \cdot 3 = 48$ Schachteln ($50 : 11$, $34 : 8$, $21 : 7$ je ganz gezählt; nicht $35\,700 : 616 \approx 58$) \quad n) $9 \cdot 3{,}5 \cdot 4{,}2 = 132{,}3$ dm³; Masse $= 132{,}3$ kg \quad o) Zylinder $\pi \cdot 18{,}5^2 \cdot 36 \approx 38\,708$ cm³; Kegel $\tfrac{1}{3} \cdot \pi \cdot 18^2 \cdot 35 \approx 11\,875$ cm³; Rest $\approx 26\,832$ cm³}
\erg{26}{a) $185 : 250 = 0{,}74 = 74\,\%$}
\erg{27}{a) $2 \cdot 6 \cdot 25 - 2 \cdot 25 = 250$ cm² (die zwei Klebeflächen fallen weg)}
\erg{28}{a) Die Klammer fehlt, die Tiefe wirkt nur auf den Halbkreis. Richtig: $(5 \cdot 4 + \tfrac{1}{2} \cdot \pi \cdot 2^2) \cdot 1{,}7 \approx 44{,}7$ m³ \quad b) Zerlegen in Quader (Wohnteil) und Dreiecksprisma (Dach), beide Volumen einzeln berechnen und addieren. \quad c) $(9 \cdot 4 + \tfrac{1}{2} \cdot \pi \cdot 2^2) \cdot 1{,}4 \approx 59{,}2$ m³ $= 59\,200$ l; $59\,200 : 1\,800 \approx 32{,}9$ h; bis Sonntag $14$ Uhr sind es $44$ h – ja \quad d) $7 \cdot 5 \cdot 3 + 5 \cdot 2 : 2 \cdot 7 = 140$ m³}
